% Wave Motions in the Ocean: Myrl's View
% Chapter 6; historical source: source/ChapmanRizzoli6.pdf.
% Source-page comments preserve printed/physical page correspondence.
% Substantive corrections are documented in ERRATA.md.

\setcounter{chapter}{5}
\setcounter{page}{149}

% -----------------------------------------------------------------------------
% Source printed page 149 / source physical page 1
% -----------------------------------------------------------------------------
\chapter{Topographic effects}

So far, we have largely ignored the effects of bottom topography. It was pointed out in
the last chapter that bottom relief appears in the shallow water vorticity equation in
the same form as the $\beta$ term and, therefore, might be expected to have similar
effects. We also introduced topography as a sudden change in depth and derived
edge-wave and Poincar\'e-wave solutions. In this chapter, we shall be more systematic
and study several types of waves which rely on variable bottom topography for their
existence. Perhaps more importantly, we shall consider the intimate relation between
variable bottom topography and stratification.

\section{Topographic Rossby waves}

We consider first a problem which was worked out by Rhines (1970) to show the combined
effects of topography and stratification. We start with the linear,

\sourcepagebreak

% -----------------------------------------------------------------------------
% Source printed page 150 / source physical page 3
% -----------------------------------------------------------------------------
hydrostatic, Boussinesq equations
\begin{align*}
 u_t-fv&=-\frac{1}{\rho_0}p_x,\\
 v_t+fu&=-\frac{1}{\rho_0}p_y,\\
 0&=-\frac{1}{\rho_0}p_z-\frac{g\rho}{\rho_0},\\
 \rho_t+w\rho_{0z}&=0,\\
 u_x+v_y+w_z&=0,
\end{align*}
where $\rho$ is the perturbation density and $\rho_0$ is considered constant except in
the density equation. The perturbation density can be eliminated to obtain
\[
 w=-\frac{1}{\rho_0N^2}p_{zt}
\]
where $N^2=-g\rho_{0z}/\rho_0$. Expressions for the velocity can be written as
\begin{align*}
 \left(\frac{\partial^2}{\partial t^2}+f^2\right)u
 &=-\frac{1}{\rho_0}p_{xt}-\frac{f}{\rho_0}p_y,\\
 \left(\frac{\partial^2}{\partial t^2}+f^2\right)v
 &=-\frac{1}{\rho_0}p_{yt}+\frac{f}{\rho_0}p_x.
\end{align*}
These expressions may be combined with continuity to yield
\[
 \left[
 p_{xx}+p_{yy}
 +\left(\frac{\partial^2}{\partial t^2}+f^2\right)
 \left(\frac{p_z}{N^2}\right)_z
 \right]_t=0.
\]
If we assume time dependence of $e^{-i\sigma t}$, then this becomes
\[
 p_{xx}+p_{yy}+(f^2-\sigma^2)
 \left(\frac{p_z}{N^2}\right)_z=0.
\]

Now consider motions confined to a channel along the $x$ axis.

\sourcepagebreak

% -----------------------------------------------------------------------------
% Source printed page 151 / source physical page 4
% -----------------------------------------------------------------------------
% Constraint-generated sloping-channel geometry.
\wavevectorart[0.70]{ch06-p151-sloping-channel}

The bottom slopes gradually across the channel with bottom slope $\alpha$. The normal
velocity must vanish at the sidewalls and at the bottom, while a rigid lid is assumed at
the surface. The boundary conditions are
\begin{wavealign}
 v=0&\quad\Rightarrow\quad i\sigma p_y+fp_x=0
 &&\text{at }y=0,L,\\
 w=0&\quad\Rightarrow\quad p_z=0
 &&\text{at }z=0,\\
 w=\alpha v&\quad\Rightarrow\quad
 i\sigma(f^2-\sigma^2)p_z
 =\alpha N^2(i\sigma p_y+fp_x)
 &&\text{at }z=-H+\alpha y.
\end{wavealign}
To proceed, we scale the variables as follows: $x,y$ by $L$; $z$ by $H$; and
$\omega=\sigma/f$. We also assume $N$ is constant. The problem then becomes
\begin{align*}
 p_{xx}+p_{yy}+\frac{1-\omega^2}{S^2}p_{zz}&=0,\\
 i\omega p_y+p_x&=0 &&\text{at }y=0,1,\\
 p_z&=0 &&\text{at }z=0,\\
 i\omega(1-\omega^2)p_z
 &=\delta S^2(i\omega p_y+p_x)
 &&\text{at }z=-1+\delta y,
\end{align*}
where $\delta=\alpha L/H$ is the scaled bottom slope and $S=NH/fL$ is the Burger
number which is a measure of the importance of stratification relative to the spatial
scales of motion. The Burger number appears in virtually all cases involving both
topography

\sourcepagebreak

% -----------------------------------------------------------------------------
% Source printed page 152 / source physical page 5
% -----------------------------------------------------------------------------
and stratification. Large $S$ means strong stratification and/or large aspect ratio
$H/L$ of the motion. Small $S$ means weak stratification and/or nearly horizontal
motions.

For the present case, we consider low frequency motions ($\omega\ll1$) over a gently
sloping bottom ($\delta\ll1$). This allows us to write the field equation and boundary
conditions as
\begin{wavealign}
 p_{xx}+p_{yy}+\frac{1}{S^2}p_{zz}&=0,\\
 p_x&=0 &&\text{at }y=0,1,\\
 p_z&=0 &&\text{at }z=0,\\
 i\omega p_z&=\delta S^2p_x &&\text{at }z=-1.
\end{wavealign}
The last boundary condition is appropriate because the fractional depth change across
the channel is small. A solution to this problem which is freely propagating in the
$x$ direction is
\[
 p=e^{ikx}\sin(n\pi y)\cosh\mu z
\]
where $\mu^2=S^2(n^2\pi^2+k^2)$ gives the vertical decay scale. Thus, strong
stratification and/or short spatial scales leads to strong bottom trapping.

% Equation-generated bottom-trapped vertical mode.
\wavevectorart[0.52]{ch06-p152-bottom-trapped-mode}

The dispersion relation is obtained by applying the bottom boundary condition
\begin{waveequation}
 \omega=-\frac{\delta kS^2}{\mu\tanh\mu}.
\end{waveequation}

\sourcepagebreak

% -----------------------------------------------------------------------------
% Source printed page 153 / source physical page 6
% -----------------------------------------------------------------------------
Notice that the waves disappear if the bottom slope vanishes $\delta\to0$ indicating
the necessity of variable topography. The phase speed is always directed so that the
waves move with the shallow water on their right in the northern hemisphere. So, for
a bottom which shoals toward the north ($+y$), the waves propagate westward ($-x$).
This is like the $\beta$-plane with nearly the same dispersion relation. If the bottom
shoals toward the south ($-y$), $\delta<0$, then the waves travel eastward ($+x$).
Thus, the effective north direction is the direction of shoaling.

The limit of weak stratification, $S\to0$, leads to $\mu\to0$ and
\begin{waveequation}
 \omega=-\frac{\delta k}{n^2\pi^2+k^2}
\end{waveequation}
or in dimensional form
\begin{waveequation}
 \sigma=-\frac{\alpha kf}
 {H\{(n\pi/L)^2+k^2\}}.
\end{waveequation}
This is the dispersion relation for Topographic Rossby waves, so named because of the
obvious similarity to planetary Rossby waves. The vertical structure, in this case,
disappears as $\mu\to0$.

\section{Bottom-trapped waves}

The waves of the previous section were indeed bottom trapped by strong stratification,
but the discussion was limited to low frequencies over a gently sloping bottom. Here
we relax these restrictions by considering waves along a sloping bottom in a
semi-infinite fluid. The motions are still assumed to be subinertial, $\sigma<f$, but
the frequency may approach $f$. This problem is also due to Rhines (1970). The field
equation for pressure is
\begin{waveequation}
 p_{xx}+p_{yy}+\frac{f^2-\sigma^2}{N^2}p_{zz}=0
\end{waveequation}

\sourcepagebreak

% -----------------------------------------------------------------------------
% Source printed page 154 / source physical page 7
% -----------------------------------------------------------------------------
where $N$ is constant. The bottom is along $z=\alpha x$ where $w=\alpha u$.

% Constraint-generated effective-slope geometry.
\wavevectorart[0.76]{ch06-p154-effective-slope}

The boundary condition along the bottom is
\begin{waveequation}
 i\sigma(f^2-\sigma^2)p_z
 =\alpha N^2(i\sigma p_x-fp_y).
\end{waveequation}
We scale $z$ by $R=N/(f^2-\sigma^2)^{1/2}$ so that $z'=zR$, and we assume a plane
wave in the $y$ direction, $e^{i\ell y}$. The problem becomes
\begin{align*}
 p_{xx}-\ell^2p+p_{z'z'}&=0,\\
 p_{z'}&=R\alpha\left(p_x-\frac{\ell f}{\sigma}p\right)
 &&\text{at }z'=R\alpha x.
\end{align*}
Thus, with stratification, $R\alpha$ is the effective bottom slope. Now
\[
 R\alpha
 =\frac{N\alpha}{(f^2-\sigma^2)^{1/2}}
 =\frac{N\alpha/f}{(1-\omega^2)^{1/2}}
 =\frac{S}{(1-\omega^2)^{1/2}}
\]
where $S=N\alpha/f$ and $\omega=\sigma/f$. Strong stratification appears as an
effectively steep bottom and vice versa. As $S\to\infty$, the bottom appears to the
motions as a vertical wall. Similarly, as $\omega\to1$, the bottom appears as a
vertical wall. The Burger number here can be thought of in the same way as in the last
section except that $H/L$ is replaced by the bottom slope $\alpha$ because there are
no distinct vertical and horizontal scales.

The angle of the bottom with respect to the horizontal is
\[
 \theta=\tan^{-1}R\alpha.
\]

\sourcepagebreak

% -----------------------------------------------------------------------------
% Source printed page 155 / source physical page 8
% -----------------------------------------------------------------------------
This allows a solution to be written as
\[
 p=e^{-i\sigma t+i\ell y
      +ik(x\cos\theta+z'\sin\theta)
      -m(z'\cos\theta-x\sin\theta)}
\]
where the time and $y$ dependences have been reinstated. In this solution, $k$ is the
wavenumber parallel to the bottom, while $m$ is the wavenumber perpendicular to the
bottom. The same solution could have been derived by first rotating the coordinate
system to be aligned with the bottom and then rotating back. Substituting this solution
into the field equation relates $k,\ell$ and $m$ as
\begin{waveequation}
 m^2=k^2+\ell^2.
\end{waveequation}
This, along with the bottom boundary condition, yields expressions for $k$ and $m$ in
terms of $\omega,\ell$ and $S$, i.e. the dispersion relation;
\begin{waveequation}
 m=\frac{\ell}{\omega}
 \left(\frac{S^2}{1-\omega^2+S^2}\right)^{1/2}
\end{waveequation}
\begin{waveequation}
 k=\frac{\ell}{\omega}
 \left[
 \frac{(S^2-\omega^2)(1-\omega^2)}
 {1-\omega^2+S^2}
 \right]^{1/2}.
\end{waveequation}
Note that for decay away from the bottom ($m>0$), $\ell$ and $\omega$ must have the
same sign. Thus, the waves propagate only in the $+y$ direction, i.e. with shoaling
water on their right just like Topographic Rossby waves. For $\omega<1$, we see that
$m$ is always real, i.e. the motions are always bottom trapped. The waves propagate
along the bottom ($k$ is real) as long as $\omega<S$. If $\omega>S$, then the waves
decay exponentially along the bottom. This means that if $S>1$, then these waves
always propagate because $\omega<1$. They become more highly bottom trapped as $S$
gets large. As $S\to0$, the waves are evanescent and less bottom trapped. As
$\omega\to0$, both $k$ and $m$ become large indicating short waves trapped close to
the bottom.

These properties suggest some interesting possibilities. Suppose a wave with frequency
$\omega<S$ propagates along the bottom and encounters a change in bottom

% -----------------------------------------------------------------------------
% Source printed page 156 / source physical page 9
% -----------------------------------------------------------------------------
\setcounter{page}{156}

slope. Without solving for the details of the solution near the corner, we know that the
wave will continue to propagate as long as the new bottom slope is such that
$\omega<S$. If $\omega>S$ on the new slope, then the wave must be reflected. Thus, we
can imagine waves being trapped on the bottom between two gently sloping regions.

% Vector reconstruction; source PDF ChapmanRizzoli6.pdf, physical page 9, trim=65bp 405bp 55bp 185bp.
\wavevectorart[0.84]{ch06-p156-bottom-slope-trapping}

This type of behavior may occur over the continental slope between the gently sloping
shelf and the gently sloping deep ocean. Of course, technically the waves would have to
be sufficiently bottom trapped so that they would not feel the surface which was
neglected in the problem. However, the surface should not fundamentally alter the wave
behavior.

\section{Continental shelf waves}

Another type of wave motion, analogous to the topographic Rossby waves but trapped at
the coast like a Kelvin wave, can occur over the continental shelf. Consider a
continental shelf which borders a flat-bottom deep ocean with depth $H$. This problem
was first considered by Buchwald and Adams (1968).

\sourcepagebreak

% -----------------------------------------------------------------------------
% Source printed page 157 / source physical page 10
% -----------------------------------------------------------------------------
% Equation-generated continental-shelf depth profile.
\wavevectorart[0.62]{ch06-p157-continental-shelf}

The $x$ axis points offshore while the $y$ axis is alongshore. The shelf-slope region has
width $L$. We start with the shallow water equations over variable topography. We
ignore stratification for now and assume that the flow is nondivergent.
\begin{align*}
 u_t-fv&=-g\eta_x,
 &v_t+fu&=-g\eta_y,\\
 (uD)_x+(vD)_y&=0.
\end{align*}
The continuity equation allows us to define a transport streamfunction as
\begin{waveequation}
 uD=\psi_y,
 \qquad
 vD=-\psi_x.
\end{waveequation}
Substituting into the momentum equations and eliminating the sea-surface displacement
yields
\[
 \left[
 \left(\frac{1}{D}\psi_x\right)_x
 +\left(\frac{1}{D}\psi_y\right)_y
 \right]_t
 +f\left[
 \left(\frac{1}{D}\right)_y\psi_x
 -\left(\frac{1}{D}\right)_x\psi_y
 \right]=0.
\]
If the topography varies only across the shelf, i.e. $D(x)$, then this becomes
\[
 \left(\psi_{xx}+\psi_{yy}-\frac{D_x}{D}\psi_x\right)_t
 +\frac{fD_x}{D}\psi_y=0.
\]
We look for plane waves propagating along the shelf, $e^{-i\sigma t+i\ell y}$, and
assume a convenient depth profile of
\[
 D=\begin{cases}
 D_0e^{2bx},&0<x<L,\\
 D_0e^{2bL},&x>L.
 \end{cases}
\]

\sourcepagebreak

% -----------------------------------------------------------------------------
% Source printed page 158 / source physical page 11
% -----------------------------------------------------------------------------
Over the shelf, the field equation reduces to
\begin{waveequation}
 \psi_{xx}-2b\psi_x
 -\left(\frac{2bf\ell}{\sigma}+\ell^2\right)\psi=0
\end{waveequation}
while in the deep ocean it becomes
\begin{waveequation}
 \psi_{xx}-\ell^2\psi=0.
\end{waveequation}
The boundary conditions are that the velocity normal to the coast must vanish and that
$\psi$ should vanish far offshore:
\begin{waveequation}
 u=0\quad\Rightarrow\quad\psi=0\quad\text{at }x=0,
\end{waveequation}
\begin{waveequation}
 \psi\to0\quad\text{as }x\to\infty.
\end{waveequation}
The solution over the shelf can be written
\[
 \psi=Ae^{b(x-L)}\sin kx.
\]
Substituting this into the $\psi$ equation yields the dispersion relation
\begin{waveequation}
 \sigma=-\frac{2bf\ell}{\ell^2+k^2+b^2}
\end{waveequation}
which looks almost identical to the Rossby wave dispersion relation showing the close
correspondence of these waves to both planetary Rossby waves and topographic Rossby
waves.

The solution in the deep sea is (since $\ell<0$)
\[
 \psi=Be^{\ell(x-L)}.
\]
The problem is closed by matching the shelf solution to the deep-sea solution. This
requires that $\psi$ and $\psi_x$ be continuous at $x=L$ which leads to
\[
 \tan kL=\frac{-k}{-\ell+b}=\frac{k}{\ell-b}.
\]

\sourcepagebreak

% -----------------------------------------------------------------------------
% Source printed page 159 / source physical page 12
% -----------------------------------------------------------------------------
This relation is satisfied at an infinite set of discrete values of $k$ for given $\ell$
and $b$. For large $k$, the roots approach $(n+\tfrac12)\pi/L$.

% Equation-generated matching-condition root plot.
\wavevectorart[0.82]{ch06-p159-shelf-root-condition}

These solutions are called \emph{continental shelf waves}. They are very much like
planetary Rossby waves and equatorial Rossby waves. Their phases all travel with the
coast on their right in the northern hemisphere. The dispersion diagram looks like

\sourceart[0.58]{ChapmanRizzoli6.pdf}{12}{135bp 265bp 120bp 350bp}

Each mode is constrained to be below some maximum frequency where the alongshore
group velocity $d\sigma_n/d\ell$ vanishes along that matched mode. Because the matching
condition makes $k=k_n(\ell)$, differentiating the dispersion relation along a branch gives
\[
 k^2+b^2-\ell^2=2\ell k\frac{dk_n}{d\ell}.
\]
The source obtains $\ell=-(k^2+b^2)^{1/2}$ by differentiating at fixed $k$; that is not
the stationarity condition along a matched modal branch and is recorded in
\texttt{ERRATA.md}.
At the maximum frequency, the alongshore group velocity is zero, meaning that energy does not
propagate even though phases still do. For waves that are longer than this wavelength
(smaller $|\ell|$), the wave energy propagates with the phase. For very long waves
$\ell\to0$, the dispersion relation becomes
\begin{waveequation}
 \sigma=-\frac{2bf\ell}{k^2+b^2}
\end{waveequation}
and the waves are nondispersive. This will be used to advantage later.

\sourcepagebreak

% -----------------------------------------------------------------------------
% Source printed page 160 / source physical page 13
% -----------------------------------------------------------------------------
Waves that are shorter than those at the frequency maxima have group velocity opposite
to the phase velocity. This means that the phases propagate forward through the group,
but the group moves in the direction with the coast on the \emph{left} in the northern
hemisphere. This is essentially identical to the result for planetary Rossby waves in
which phase always propagates to the west, but the group velocity may be westward or
eastward depending on the wavelength of the Rossby wave. One difference is that
continental shelf waves have discrete cross-shelf modes whereas unbounded Rossby waves
form a continuum. Of course, Rossby waves would be discretized if they were constrained
to a channel of some sort. The coast acts as this sort of constraint for continental
shelf waves. Notice that the frequency for each mode approaches zero as the waves
become very short, i.e. $\sigma\to0$ as $\ell\to-\infty$.

We have made a special choice for the bottom topography which made the problem rather
simple by giving constant coefficients to the equation for $\psi$. It can be shown that
the present results are but a special case of the results for the more general divergent
equations with arbitrary cross-shelf topography. The equations are
\begin{align*}
 u_t-fv&=-g\eta_x,
 &v_t+fu&=-g\eta_y,\\
 \eta_t+(uD)_x+(vD)_y&=0.
\end{align*}
If we assume that the topography does not vary along the shelf, i.e.
$\partial D/\partial y=0$, and look for plane waves of the form
$e^{-i\sigma t+i\ell y}$, then the problem becomes
\[
 (D\eta_x)_x-K\eta=0,
 \qquad
 K=\frac{f\ell}{\sigma}D_x+\ell^2D+\frac{f^2-\sigma^2}{g}.
\]
The boundary conditions are
\begin{waveequation}
 uD=0\quad\Rightarrow\quad
 D\left(\eta_x-\frac{f\ell}{\sigma}\eta\right)=0
 \quad\text{at }x=0,
\end{waveequation}
\begin{waveequation}
 \eta\to0\quad\text{as }x\to\infty.
\end{waveequation}

\sourcepagebreak

% -----------------------------------------------------------------------------
% Source printed page 161 / source physical page 15
% -----------------------------------------------------------------------------
which represent no flow through the coast and coastal trapping. Huthnance (1975) has
shown that, provided $D$ increases monotonically offshore, this eigenvalue problem
yields an infinite discrete set of continental shelf waves which have the same general
properties as those for the special case above. Further, exactly one Kelvin wave exists
which can propagate at both sub- and super-inertial frequencies. Also, there is an
infinite discrete set of edge waves, all at super-inertial frequencies, which can
propagate in either direction. They occur outside a continuum of Poincar\'e waves. The
complete dispersion diagram looks like

\sourceart[0.70]{ChapmanRizzoli6.pdf}{15}{105bp 325bp 85bp 225bp}

Notice the obvious similarity to the dispersion diagram for equatorial waves. In fact,
most of the waves in the equatorial dispersion diagram have counterparts along the
coast, except that there is no Yanai wave along a coast. Thus, to many researchers, the
coastal region is essentially the same as the equator, but turned sideways. There is
another important distinction, however, which we shall discuss next. That is the role
of stratification. Our results from the equator were easily generalizable to a stratified
ocean because the bottom was flat, so we could make use of the expansion in vertical
modes and simply use a different equivalent depth to study higher modes. In contrast,
waves trapped at the coast depend on the variations in topography to exist. This, along
with the intimate relationship between topography and stratification which we

\sourcepagebreak

% -----------------------------------------------------------------------------
% Source printed page 162 / source physical page 16
% -----------------------------------------------------------------------------
discussed in the previous two sections, suggests that the inclusion of stratification may
not be trivial for continental shelf waves.

\section{Coastal-trapped waves}

In order to add stratification to the continental shelf wave problem, we must return to
the linear, hydrostatic, Boussinesq equations with which we started the chapter. We
consider a coastline oriented along the $y$ axis with $x$ pointing offshore.

% Vector reconstruction; source PDF ChapmanRizzoli6.pdf, physical page 16, trim=155bp 360bp 125bp 270bp.
\wavevectorart[0.56]{ch06-p162-coastal-geometry}

We assume that the topography varies only across the shelf and look for free waves
propagating in $y$, i.e. $e^{-i\sigma t+i\ell y}$. The equation and boundary conditions
in terms of pressure are
\begin{waveequation}
 p_{xx}-\ell^2p+\frac{f^2-\sigma^2}{N^2}p_{zz}=0,
\end{waveequation}
\begin{waveequation}
 (f^2-\sigma^2)p_z
 +N^2D_x\left(p_x-\frac{f\ell}{\sigma}p\right)=0
 \quad\text{at }z=-D(x),
\end{waveequation}
\begin{waveequation}
 p_z=0\quad\text{at }z=0,
\end{waveequation}
\begin{waveequation}
 p\to0\quad\text{as }x\to\infty.
\end{waveequation}
We have taken $N$ to be constant and applied a rigid lid. However, all of the following
analysis can be generalized to the case of variable $N$ and a free surface (see
Huthnance, 1978).

% -----------------------------------------------------------------------------
% Source printed page 163 / source physical page 17
% -----------------------------------------------------------------------------
\setcounter{page}{163}

We scale the variables as follows: $x,y$ by $L$; $z,D$ by $H$; and
$\omega=\sigma/f$. The equations become
\begin{align*}
 p_{xx}-\ell^2p+\frac{1-\omega^2}{S^2}p_{zz}&=0,\\
 (1-\omega^2)p_z
 +S^2D_x\left(p_x-\frac{\ell}{\omega}p\right)&=0
 &&\text{at }z=-D(x),\\
 p_z&=0 &&\text{at }z=0,\\
 p&\to0 &&\text{as }x\to\infty,
\end{align*}
where $S=NH/fL$ as before. For general $D(x)$, this eigenvalue problem must be solved
numerically. In fact, only a couple of special cases of $D$ are known which give
analytical solutions. And these are rather unusual in their properties, so we will not
study them. However, a number of important features of the free-wave solutions can be
determined without solving the complete problem. These are all due to Huthnance
(1978).

\begin{enumerate}
 \item There is a singly infinite discrete set of wave modes for any choice of topography
       and stratification. These are called \emph{coastal-trapped waves}.
 \item Increased stratification, all else being equal, increases the wave frequency and
       makes the wave structure more horizontal.
 \item The dispersion curves for all modes approach the same frequency as the
       wavelength decreases. This frequency is given by
       \[
       \lim_{\ell\to-\infty}\omega=S\max[D_x].
       \]
 \item The short waves (large $\ell$) are identical to the bottom-trapped waves found
       by Rhines (1970).
\end{enumerate}

These results give the following dispersion diagram for the general case

\sourcepagebreak

% -----------------------------------------------------------------------------
% Source printed page 164 / source physical page 18
% -----------------------------------------------------------------------------
% Vector reconstruction; source PDF ChapmanRizzoli6.pdf, physical page 18, trim=75bp 435bp 65bp 70bp.
\wavevectorart[0.82]{ch06-p164-ctw-dispersion-family}

The second and third results say that the dispersion curves will go higher and higher
with increasing stratification, and if $S\max[D_x]>1$ then all dispersion curves go to
the inertial frequency, $\omega=1$. In dimensional form, this is
$(N/f)\max[D_x]$ where $D_x$ is the actual (not scaled) bottom slope.

% Vector reconstruction; source PDF ChapmanRizzoli6.pdf, physical page 18, trim=110bp 240bp 95bp 385bp.
\wavevectorart[0.68]{ch06-p164-stratification-effect}

This has profound effects on the nature of the waves. They are no longer restricted to
be below a maximum frequency. Now they may occur at any subinertial frequency, but
they are limited in length by the wavenumber at which the dispersion curve reaches
$f$. That is, each mode must be longer than a certain length to be a free wave. Consider
the change that this makes on a scattering problem.

\sourcepagebreak

% -----------------------------------------------------------------------------
% Source printed page 165 / source physical page 19
% -----------------------------------------------------------------------------
% Vector reconstruction; source PDF ChapmanRizzoli6.pdf, physical page 19, trim=70bp 510bp 65bp 45bp.
\wavevectorart[0.82]{ch06-p165-scattering-by-stratification}

If the stratification is weak, then waves may exist which propagate energy in either
direction because the group velocity changes sign. Thus, energy may be reflected as
well as transmitted. If the stratification is strong so that all of the dispersion curves
go to $f$, then the energy can only propagate in one direction. No energy can be
reflected from the topography, no matter how tortuous the topographic variations. It
turns out that in the ocean, $ND_x/f$ is often order 1, especially at low latitudes where
$f$ is small. Typically, at mid-latitudes, $N/f$ is order 10 to 100, while $D_x$ is order
$10^{-3}$ over the shelf but more like 0.02--0.04 over the continental slope. Remember
that $N/f$ times the maximum of $D_x$ is the important value.

Before leaving this problem it is useful to consider two limiting cases.

\noindent\textbf{Case A: $S\to0$.} If $S$ is small, we can expand the solution in
powers of $S^2$ as follows
\[
 p(x,z)=p_0(x)+S^2p_1(x,z)+O(S^4).
\]
Substituting into the full equations produces
\[
 O(1):\qquad p_{0zz}=0
\]
with $p_{0z}=0$ at both $z=0$ and $z=-D$. This means that $p_{0z}=0$ everywhere,
i.e. the solution $p_0$ is vertically uniform. At the next order, we have
\[
 O(S^2):\qquad p_{0xx}-\ell^2p_0+(1-\omega^2)p_{1zz}=0.
\]

\sourcepagebreak

% -----------------------------------------------------------------------------
% Source printed page 166 / source physical page 20
% -----------------------------------------------------------------------------
\[
 (1-\omega^2)p_{1z}
 +D_x\left(p_{0x}-\frac{\ell}{\omega}p_0\right)=0
 \quad\text{at }z=-D,
\]
\[
 p_{1z}=0\quad\text{at }z=0.
\]
The field equation can be integrated in $z$, since $p_0$ is independent of $z$, and
combined with the surface and bottom boundary conditions to yield
\begin{waveequation}
 (Dp_{0x})_x-\left(\frac{\ell}{\omega}D_x+\ell^2D\right)p_0=0
\end{waveequation}
which is precisely the same equation that was derived for continental shelf waves, but
now with a rigid lid. Thus, as we would expect, the stratified problem reduces to the
barotropic problem in the limit of weak stratification.

\noindent\textbf{Case B: $S\to\infty$.} Based on our previous experience with
stratification effects, we expect strong stratification to lead to strong bottom-trapping,
i.e. short vertical scales. Basically this occurs because the stratification inhibits
vertical motions. Therefore, we make a change of variables to
\[
 \xi=x-D^{-1}(-z),
 \qquad
 \eta=Sz
\]
where $D^{-1}$ is the inverse of the depth function. The new variable $\xi$ represents
the horizontal distance from the bottom. The equations become, for large $S$,
\begin{align*}
 p_{\xi\xi}+(1-\omega^2)p_{\eta\eta}-\ell^2p&=0,\\
 D_x\left(p_\xi-\frac{\ell}{\omega}p\right)&=0 &&\text{at }\xi=0,\\
 p_\eta&=0 &&\text{at }\eta=0.
\end{align*}
A solution which satisfies the first two and the requirement that $p_z=0$ on the deep
ocean bottom at $z=-1$ is
\[
 p=e^{\ell\xi/\omega}
 \cos\left[\frac{\ell}{\omega}(\eta+S)\right].
\]

\sourcepagebreak

% -----------------------------------------------------------------------------
% Source printed page 167 / source physical page 21
% -----------------------------------------------------------------------------
Remember that $\ell<0$ in the present formulation, so the solution does decay offshore.
The surface boundary condition provides the dispersion relation of
\begin{waveequation}
 \boxed{\omega=-\frac{S\ell}{n\pi}},
 \qquad n=1,2,\ldots
\end{waveequation}
where the minus sign is the positive-frequency form for the source's $\ell<0$ convention;
the source prints the same relation without the minus sign. See \texttt{ERRATA.md}.
This is precisely the same dispersion relation as for baroclinic Kelvin waves with
constant $N$.

Thus, each coastal-trapped wave mode behaves like a continental shelf wave when the
stratification is weak, and then passes smoothly to a baroclinic Kelvin wave when the
stratification is strong.

% Constraint-generated weak-to-strong stratification mode transition.
\wavevectorart[0.68]{ch06-p167-mode-transition}

Consider a free wave travelling north along an eastern ocean boundary with constant
$N$ and uniform $D(x)$. At low latitudes, the wave looks like a baroclinic Kelvin wave
because $S$ is large. However, as the wave moves north, $S$ decreases and the wave
looks more and more like a continental shelf wave. Of course, we have neglected the
$\beta$ effect which would change the entire problem. So, we cannot take our thought
experiment too far.

\section{Wind-forced, long waves}

We have not discussed how shelf or coastal-trapped waves might be generated. Over the
past 15 years or so, a very elegant theory has evolved which suggests that the
alongshelf component of the surface wind stress is an important driving mechanism.

\sourcepagebreak

% -----------------------------------------------------------------------------
% Source printed page 168 / source physical page 22
% -----------------------------------------------------------------------------
This theory has proven quite successful in predicting shelf currents, so we will examine
the basics of it. We will consider only the simplest form of the theory by Gill and
Schumann (1974), but you should keep in mind that it has been generalized to a much
more realistic setting.

We consider a barotropic (homogeneous) ocean as we did for the continental shelf waves.

% Vector reconstruction of the long-wave wind-forcing geometry.
\wavevectorart[0.58]{ch06-p168-wind-forced-shelf}

In addition, we assume that the motions occur at frequencies much less than the inertial
frequency, i.e. $\sigma\ll f$, and that the alongshelf variations occur on a much larger
scale than the cross-shelf motions, i.e.
$\partial/\partial y\ll\partial/\partial x$. These assumptions constitute the
\emph{long-wave approximation}. In terms of the free-wave dispersion diagram, we are
assuming that the waves are at small $\sigma$ and small $\ell$. The equations of motion
are
\begin{wavealign}
 -fv&=-g\eta_x,\\
 v_t+fu&=-g\eta_y+\frac{\tau^y}{D},\\
 (uD)_x+(vD)_y&=0.
\end{wavealign}
We have also assumed a rigid lid, and imposed an alongshelf wind stress $\tau^y$.
Notice that the long-wave approximation has rendered the alongshelf flow to be in
geostrophic balance. This turns out to be a good approximation over most continental
shelves.

We define a streamfunction as
\[
 uD=\psi_y,
 \qquad
 vD=-\psi_x.
\]

\sourcepagebreak

% -----------------------------------------------------------------------------
% Source printed page 169 / source physical page 23
% -----------------------------------------------------------------------------
which results in an equation for $\psi$ of (with $D_y=0$)
\[
 \left(\frac{\psi_x}{D}\right)_{xt}
 +\frac{fD_x}{D^2}\psi_y
 =\frac{D_x}{D^2}\tau^y.
\]
The boundary conditions are that $\psi=0$ at $x=0$, i.e. no flow through the coast,
and $\psi_x=0$ at $x=L$ which comes from matching $\psi_x$ at $x=L$. The $x$ and
$y$ length scales are both order $\ell^{-1}$ in the deep ocean, so
$\psi_x\sim\ell\psi\simeq0$ at $x=L$.

To solve this problem, we first look at free-wave solutions, i.e. $\tau^y=0$. Then we
separate variables by writing
\[
 \psi(x,y,t)=\phi(y,t)F(x).
\]
The field equation becomes
\begin{waveequation}
 \phi_t\left(\frac{F_x}{D}\right)_x
 +\frac{fD_x}{D^2}\phi_yF=0
\end{waveequation}
for which the separation works only if
\begin{waveequation}
 \frac{1}{c}\phi_t-\phi_y=0
\end{waveequation}
\begin{waveequation}
 \left(\frac{F_x}{D}\right)_x
 +\frac{fD_x}{cD^2}F=0
\end{waveequation}
where $c$ is a separation constant. The boundary conditions become $F=0$ at $x=0$
and $F_x=0$ at $x=L$.

The problem for $F$ is a Sturm--Liouville eigenvalue problem and it can be shown that
the eigenfunctions, $F_n$, satisfy the orthogonality relation
\begin{waveequation}
 \int_0^L\frac{D_x}{D^2}F_nF_m\,dx=\delta_{nm}
\end{waveequation}
where $\delta_{nm}$ is the Kronecker delta. Each $F_n$ corresponds to the cross-shelf
structure of a free-wave mode. The problem for $\phi$ is just a first-order wave equation
with the solution
\begin{waveequation}
 \phi=\phi_0(y+ct).
\end{waveequation}

\sourcepagebreak

% -----------------------------------------------------------------------------
% Source printed page 170 / source physical page 24
% -----------------------------------------------------------------------------
where $\phi_0$ is any function. Thus, we see that each wave mode need not be sinusoidal
in shape and that the eigenvalue $c$ is simply the phase speed of the free wave. The
waves move in the $-y$ direction as expected and they are nondispersive (as shown for
the small $\ell$ limit of the continental shelf wave problem).

These results allow the forced problem to be solved by expanding $\psi$ in the set of
free-wave modes
\[
 \psi(x,y,t)=\sum_n\phi_n(y,t)F_n(x).
\]
Substituting this into the field equation and using the orthogonality condition of the
free modes yields
\begin{waveequation}
 \frac{1}{c_n}\phi_{nt}-\phi_{ny}=-b_n\tau^y
\end{waveequation}
where
\begin{waveequation}
 b_n=\frac{1}{f}\int_0^L\frac{D_x}{D^2}F_n\,dx
\end{waveequation}
are the wind-coupling coefficients which tell how well the wind stress drives each mode.
The solution to this equation is
\[
 \phi_n(y,t)
 =\phi_n\left(0,t+\frac{y}{c_n}\right)
 +b_n\int_0^y
 \tau^y\left(\xi,t+\frac{y-\xi}{c_n}\right)d\xi.
\]
This solution simply says that $\phi(y,t)$ is given by the $\phi$ that propagated into
the domain at the origin of integration plus the integrated effect of the wind generating
free waves along the coast. Thus, the entire wind-forced problem has boiled down to a
rather simple integral of the wind stress displaced in time by the period needed for the
free wave to propagate from its generation location $\xi$ to the prediction site $y$.

The overall solution procedure is as follows. The free-wave phase speeds and cross-shelf
structures are computed from the eigenvalue problem. These are used to compute the
$b_n$. Then the first-order wave equation is integrated for each mode and

\sourcepagebreak

% -----------------------------------------------------------------------------
% Source printed page 171 / source physical page 25
% -----------------------------------------------------------------------------
the streamfunction is reconstructed as the appropriate summation of modes. Crucial to
this approach is the long-wave approximation which renders the waves nondispersive
allowing the separation of variables. This approach does not work for dispersive waves.
This theory has been extended to a remarkable degree of sophistication in which
alongshelf variations in stratification, bottom topography and bottom friction have been
incorporated.